\chapter{INTRODUCTION}
\pagenumbering{arabic}

\section{Background Introduction}

Sign languages are complete natural languages with their own vocabulary and grammar. They are not signed versions of spoken languages, and they do not translate word for word into them. Nepali Sign Language (NSL) is the language of the deaf community in Nepal and differs substantially from American, British, or Indian sign languages. It grew within Nepal's deaf community after the first deaf school opened in Kathmandu in the 1960s, and the standard gesture set is published by the National Federation of the Deaf Nepal \cite{acharya2003nsldict, sunuwar2024nsl23}.

\noindent The community that uses it is larger than official figures suggest. The 2011 national census recorded roughly 80{,}000 deaf or hard-of-hearing Nepalis. The National Federation of the Deaf Nepal puts the real figure above 300{,}000 and attributes the difference to the undercounting of a condition that is not visible to an enumerator \cite{nepalitimes2022, hi2023}. For most of these people sign language is the primary means of communication, and almost none of the hearing population can understand it.

\noindent This gap is inconvenient in ordinary life and dangerous in an emergency. A deaf person who needs to report a fire, ask for an ambulance, or tell a bystander that they cannot breathe has no reliable way to do so if no interpreter is present, and interpreters are scarce and cannot be summoned on demand. The time spent failing to communicate is time that matters.

\noindent Recognition technology has advanced quickly for well-resourced sign languages, particularly American Sign Language, but NSL has received very little of that attention. Only one public NSL dataset exists and it covers the fingerspelling alphabet \cite{sunuwar2024nsl23}. The handful of published NSL recognition systems target individual letters or static handshapes \cite{mali2018nsl, shrestha2024nsl, poudel2025nsl}. Nothing published addresses emergency communication, and nothing addresses whole phrases, which are what a person actually needs to sign under pressure.

\noindent Two developments make a lightweight system practical. Pretrained hand and body tracking such as MediaPipe \cite{zhang2020mediapipe} extracts the geometry of a signing person from an ordinary camera without gloves or depth sensors, and reduces each video frame to a few hundred numbers. Sequence models such as the bidirectional Long Short-Term Memory network \cite{hochreiter1997lstm, graves2005bilstm} classify the resulting motion directly. Together they allow a recognition system to be trained from a few hundred recordings and to run on hardware people already own.

\noindent This project applies that combination to a narrow, high-value problem: recognising six NSL emergency phrases from a live camera feed and delivering each as on-screen text and speech. The six phrases are \textit{I can't breathe} (\nepali{म सास फेर्न सक्दिन}), \textit{The building is on fire} (\nepali{घरमा आगो लाग्यो}), \textit{Call the police} (\nepali{प्रहरीलाई बोलाउनुहोस्}), \textit{I need an ambulance} (\nepali{मलाई एम्बुलेन्स चाहियो}), \textit{Help me / I am in danger} (\nepali{मलाई मद्दत गर्नुहोस्}), and \textit{I need to go to the toilet} (\nepali{मलाई शौचालय जानु छ}). They span medical, fire, crime, and basic-need situations. Their signed form was established with Deaf experts before any data was recorded, so what the system recognises is real NSL and not an invention of the authors.

\section{Problem Statement}

In an emergency, a deaf person in Nepal has no reliable way to make an urgent need understood. Bystanders and first responders do not understand Nepali Sign Language, and an interpreter is almost never present when an emergency begins. At the same time, no recognition system or dataset exists for NSL emergency phrases; published work stops at the fingerspelling alphabet and static letters \cite{sunuwar2024nsl23, shrestha2024nsl, poudel2025nsl}.

\noindent What is missing is a low-cost tool that recognises essential NSL emergency phrases as they are signed and presents them in a form a hearing person understands immediately.

\noindent This project was necessary because nothing in the published record supplies that tool or the means to build one. There is no NSL emergency-phrase dataset to train on, no NSL system that recognises motion-based phrases rather than static letters, and no NSL result reported in a way that establishes how a system behaves on a person it has never seen. A tool that announces emergencies also has to be able to stay silent when the input is not a sign at all, and no published NSL system addresses that either. Section~\ref{sec:researchgap} states this gap against the reviewed literature and Section~\ref{sec:contributions} states what this project adds to close it.

\section{Objectives}
\label{sec:objectives}

To design and implement a system that establishes the signed form of six NSL emergency phrases with NSL experts, builds a custom landmark-based dataset from multiple signers, and trains and evaluates a BiLSTM classifier on signers absent from training, deployed in a mobile application that recognises phrases on-device and outputs them as text and speech.

\section{Main Contributions}
\label{sec:contributions}

This project makes four contributions. The first is the principal one, and the one against which the rest of this report should be read.

\begin{enumerate}
    \item \textbf{Identity-linked, signer-independent evaluation of NSL recognition.} Every clip in the dataset is tagged with the identity of the person who recorded it. The classifier is then trained and scored eight times over, and each of those models is tested only on the one signer absent from its own training data. This gives 95.18\% mean accuracy across the eight folds and 95.36\% pooled over all 840 predictions. Trained on the same data under a random split, the same architecture reports 97.62\%. That difference of 2.3 percentage points is a direct measurement of how much a signer's identity leaks into a headline accuracy when clips from one person sit on both sides of the split. No published NSL recognition work separates signers between training and test, so no published NSL accuracy establishes what a system does for a new user. This report is the first to make that measurement for NSL, and it is the contribution most worth carrying into further work.

    \item \textbf{The first landmark-based NSL emergency-phrase dataset.} 840 clips from eight signers across seven classes, balanced at exactly 120 clips per class. The signed form of every phrase was fixed in person with Deaf experts before any recording began, so the dataset records real NSL rather than an invention of the authors. No public NSL dataset of emergency phrases existed at the start of this work.

    \item \textbf{Explicit rejection of input that is not a target phrase.} A seventh class trained on non-signing motion, combined with a three-outcome decision rule, lets the system report an input as unrecognised instead of naming an emergency phrase for it. For a tool whose output is a spoken emergency announcement, the ability to stay silent matters as much as the ability to recognise.

    \item \textbf{A recogniser that runs entirely on an ordinary phone.} The trained network is exported to TensorFlow Lite and shipped inside an Android application that extracts landmarks, classifies them, displays the phrase in Nepali and English, and speaks it aloud in Nepali, with no network connection at any point.
\end{enumerate}

\section{Scope of the Project}

The system recognises exactly seven classes: the six emergency phrases listed above and a seventh negative class covering non-signing motion, idle movement, and gestures outside the target set. Each input is one complete phrase performed once. Recognition works from an ordinary camera with no gloves, markers, or depth sensors, and the model runs without a GPU.

\noindent The seventh class deserves a note. A classifier trained on six phrases will assign one of those six to any input it receives, including someone scratching their head. For a system whose output is an emergency announcement, that behaviour is unacceptable. The negative class gives the model explicit examples of what is not a sign, so the confidence threshold described in Chapter~\ref{ch:methodology} can withhold output for inputs unlike anything in the six target phrases.

\noindent Outside the scope are sentence-level or conversational translation, the full NSL vocabulary, recognition of more than one person at a time, and facial expression. Facial landmarks were excluded after the NSL consultants confirmed that none of the six target phrases depends on facial expression to be understood. The system is a proof of concept for a fixed phrase set and is built so that phrases, signers, and deployment targets can be added later.
